\documentclass[11pt]{article}
\usepackage[a4paper,margin=27mm]{geometry}
\usepackage{amsmath,amssymb,amsthm,booktabs,array}
\usepackage[colorlinks=true,linkcolor=blue,citecolor=blue,urlcolor=blue]{hyperref}
\usepackage[T1]{fontenc}
\usepackage{lmodern}

\newtheorem{theorem}{Theorem}
\newtheorem{lemma}{Lemma}
\newtheorem{proposition}{Proposition}
\newcommand{\CNOT}{\mathrm{CNOT}}
\newcommand{\supp}{\operatorname{supp}}
\newcommand{\tr}{\operatorname{tr}}
\newcommand{\Span}{\operatorname{span}}
\newcommand{\Frob}{\mathrm F}

\title{Local Observables Obstruct Short Diagonalizers\\of the Four-Qubit Cycle}
\author{Standalone research draft}
\date{\today}

\begin{document}
\maketitle

\begin{abstract}
How many entangling gates are needed to turn the eigenbasis of a cyclic
qubit permutation into a computational basis? For four qubits, we show
that an exact ancilla-free diagonalizer needs at least six CNOTs, even
with arbitrary one-qubit gates and unrestricted connectivity. The idea
is to follow a suitably chosen local input observable. Diagonalization
restricts the spectral sectors through which this observable can pass,
whereas a short circuit restricts its physical support. These constraints
are incompatible for every circuit with at most five CNOTs. The proof
uses elementary operator identities together with finite, exact algebraic
certificates and an exhaustive classification of interaction schedules.
It establishes a lower bound; the minimum CNOT count remains open.
\end{abstract}

\section{The diagonalization problem}

The cyclic shift is easy to specify in the computational basis, but its
eigenvectors need not be easy to prepare. We ask for the entangling cost
of changing to an eigenbasis, allowing the circuit to choose both the
ordering of the eigenvalues and the basis within each degenerate sector.
This freedom matters: a lower bound for one prescribed eigenbasis would
not settle the diagonalization problem.

Label the qubits $0,1,2,3$, with the leftmost bit most significant, and let
\begin{equation}\label{eq:shift}
V\lvert x_0x_1x_2x_3\rangle
 =\lvert x_3x_0x_1x_2\rangle .
\end{equation}
Thus $V^4=I$.  A diagonalizer is a four-qubit unitary $T$ for which
$TVT^\dagger$ is diagonal in the computational basis.  We count CNOTs in an
exact, ancilla-free circuit for $T$; between CNOTs we permit arbitrary
one-qubit unitaries, and each CNOT can have either direction on any physical
pair.  There are no measurements, postselection, or implicit wire
permutations.  The diagonal entries may be ordered arbitrarily, and no basis
is prescribed within an eigenspace of $V$.

\begin{theorem}\label{thm:main}
Let $T$ be any such exact diagonalizer of $V$.  If a circuit for $T$ contains
$m$ CNOTs, then $m\geq6$.
\end{theorem}

Equivalently, if $c_4$ is the minimum in this gate model, then
$c_4\geq6$.  The supplied cyclic-shift draft proves the general bound
$c_n\geq n-1$, which gives three at $n=4$ \cite[Theorem~2]{cycle-draft}.
The proof below is a stronger four-qubit statement.  A checked construction
with fourteen CNOTs is described by the companion gate list
\path{topology14_exact_matchgate_rational.json} and checked over
$\mathbb Q(\zeta_{96})$ by \path{exact_check.py}; hence
$6\leq c_4\leq14$.  These claims concern diagonalization of $V$, not a
simultaneous diagonalization of total spin.

For the proof it is convenient to set $U=T^\dagger$.  Inverting a circuit
preserves its CNOT count, and
\begin{equation}\label{eq:diagonal}
U^\dagger VU=D
\end{equation}
for some diagonal $D$.  We read the CNOTs of $U$ in chronological order.
Let $P_a$ be the spectral projector of $V$ for eigenvalue $i^a$,
$a\in\{0,1,2,3\}$.  Write
\begin{equation}\label{eq:projector}
M_a=4P_a=\sum_{r=0}^{3}i^{-ar}V^r.
\end{equation}

\section{What a local input observable remembers}

The circuit $U$ maps computational basis states to eigenvectors of $V$.
Instead of following all sixteen states, we follow just one observable
on one input wire. A one-qubit operator connects only two computational
basis strings at a time. This elementary observation survives the change
of basis and gives the spectral constraint below.

An \emph{input Pauli axis} on wire $j$ is
$p_j=\mathbf n\cdot\boldsymbol\sigma_j$, where
$\mathbf n\in\mathbb R^3$ has unit norm, and $\boldsymbol\sigma=(X,Y,Z)$.
We write $\sigma(\mathbf n)=\mathbf n\cdot\boldsymbol\sigma$. Its physical image is
$A=Up_jU^\dagger$.  It is a traceless Hermitian involution.

\begin{lemma}[Spectral-path identity]\label{lem:path}
For every input Pauli axis $p_j$ and pairwise distinct $a,b,c$,
\begin{equation}\label{eq:path}
P_a A P_b A P_c=0.
\end{equation}
No nonzero Pauli axis supported on one physical qubit satisfies all these
identities.
\end{lemma}

\begin{proof}
Conjugate the first identity by $U^\dagger$.  Each $U^\dagger P_aU$ is a
diagonal spectral projector of $D$.  An input Pauli axis connects a basis
string only to itself or to the string obtained by flipping bit $j$.
Traversing two distinct spectral labels requires the flip twice, so the
third label equals the first.  This contradicts pairwise distinctness.

For the second assertion, consider $Z_0$ on the one-excitation subspace.
Let $e_r$ have its excitation at wire $r$, and set
\begin{equation}\label{eq:fourier}
\phi_a=\frac12\sum_{r=0}^{3}i^{-ar}e_r.
\end{equation}
Then $V\phi_a=i^a\phi_a$ and, for $a\ne b$,
$\langle\phi_a,Z_0\phi_b\rangle=-1/2$.  As $Z_0$ preserves excitation
number, the matrix element of $P_0Z_0P_1Z_0P_2$ between $\phi_0$ and
$\phi_2$ is $1/4$.  It is nonzero.  A cyclic translation takes this result
to any physical wire.  A collective one-qubit rotation $R^{\otimes4}$
commutes with $V$ and takes $Z$ to any Pauli axis, proving the claim.
\end{proof}

The last part of Lemma~\ref{lem:path} is a useful warm-up: an observable
that stays local cannot be the image of a local input axis. The next
statement gives a different restriction, now using the particular input
axis $Z_j$. Since $Z_j$ commutes with $D$, its image must respect the
physical cyclic symmetry.

\begin{lemma}[Full support of a commuting observable]\label{lem:support}
Every non-scalar physical operator commuting with $V$ has support on all
four wires.  In particular, $UZ_jU^\dagger$ has full support for every
input wire $j$.
\end{lemma}

\begin{proof}
The minimal physical support of an operator is unique.  Conjugation by
$V$ cyclically translates that support.  If the operator commutes with
$V$, its support must be invariant under this transitive four-cycle; hence
it is empty or all four wires.  Empty support means a scalar.  Finally,
$Z_j$ commutes with diagonal $D$, so $UZ_jU^\dagger$ commutes with $V$ by
Eq.~\eqref{eq:diagonal}, and it is not scalar.
\end{proof}

\begin{lemma}[Minimum incidence]\label{lem:degree}
Every input wire meets at least two CNOTs in a circuit for $U$.
Moreover, the forward support cone of each input wire, formed by adding
both endpoints whenever its current support touches a CNOT pair, reaches
all four wires.
\end{lemma}

\begin{proof}
If wire $j$ meets no CNOT, an input Pauli axis stays one-body, contrary
to Lemma~\ref{lem:path}.  If it meets exactly one CNOT, choose the input
axis so that immediately before this gate it is $Z$ on a control or $X$
on a target.  It commutes through the gate and remains one-body, giving
the same contradiction.  One-qubit gates before or after the CNOT merely
rotate the axis.  The support-cone claim follows from
Lemma~\ref{lem:support}: a CNOT can enlarge physical support only across
its incident pair.  The stated symmetric update is an overapproximation
that permits either CNOT direction.
\end{proof}

Two incidences per wire already require four CNOTs. To exclude four and
five, we choose an axis that commutes through its first incident CNOT.
Only the later gates can then spread it. A surviving two-body image on
an opposite pair needs one further test: the half-shift $V^2$, which
swaps $(0,2)$ and $(1,3)$, detects such an image through its trace.

\begin{lemma}[Half-shift trace obstruction]\label{lem:trace}
Suppose $A=Up_jU^\dagger$ has proper physical support and commutes with
$V^2$.  Then $\tr(V^2A)=0$.  For opposite physical wires $r,s$ and any
unit Bloch vector $\mathbf n$,
\begin{equation}\label{eq:trace-product}
\tr\!\left(V^2\sigma(\mathbf n)_r\sigma(\mathbf n)_s\right)=4.
\end{equation}
Consequently $A$ cannot equal
$\pm\sigma(\mathbf n)_r\sigma(\mathbf n)_s$ on an opposite pair.
\end{lemma}

\begin{proof}
If $p_j=\pm Z_j$, then $A$ commutes with $V$ by
Eq.~\eqref{eq:diagonal}, contradicting its proper support and
Lemma~\ref{lem:support}.  Thus $p_j$ has a nonzero off-diagonal entry
between the two states of bit $j$.  The commutator $[A,V^2]=0$ implies
$[p_j,D^2]=0$.  Since $D^2$ is diagonal, its entries agree across every
edge of the computational hypercube that flips bit $j$.  The diagonal
entries of $p_j$ on the two ends of such an edge are opposite; pairing
terms in $\tr(D^2p_j)$ gives
\begin{equation}\label{eq:trace-zero}
\tr(V^2A)=\tr(D^2p_j)=0.
\end{equation}
For Eq.~\eqref{eq:trace-product}, write $V^2$ as the product of SWAPs
on the two opposite pairs.  The identity
$\tr(\mathrm{SWAP}\, B\otimes C)=\tr(BC)$ gives $2$ for the pair
carrying the two equal Pauli axes and $2$ for the identity pair.
\end{proof}

\section{The images left by one or two later CNOTs}

We now translate the circuit restriction into small spaces of operators.
The spaces below deliberately contain more operators than an actual
circuit can produce. Thus their exclusion covers arbitrary local gates,
without choosing angles or fitting a numerical ansatz.

The two propositions in this section are the computer-assisted part of
the argument. Their precise finite matrix and polynomial tests are in
Appendix~\ref{app:gram}. Here we state the consequences needed for the
circuit proof. All spans are over $\mathbb R$, and $\sigma_\mu$ ranges
over $X,Y,Z$.

\begin{proposition}[One-CNOT image]\label{prop:terminal}
Suppose a selected input Pauli axis commutes through its first incident
CNOT and encounters exactly one later incident CNOT, on physical pair $(j,k)$.
Assume no subsequent CNOT touches either $j$ or $k$. Up to a collective one-qubit rotation, its final image belongs to
\begin{equation}\label{eq:pair-space}
\mathcal L_{jk}=\Span_{\mathbb R}
 \{\sigma_\mu^{(j)},\ \sigma_\mu^{(j)}Z_k:
   \mu\in\{X,Y,Z\}\}.
\end{equation}
If $(j,k)$ is adjacent on the physical cycle, no nonzero member of this
space satisfies Eq.~\eqref{eq:path}.  If $(j,k)$ is opposite, every
nonzero member satisfying Eq.~\eqref{eq:path} is, after undoing the
collective rotation, proportional to a same-axis product
$\sigma(\mathbf n)_j\sigma(\mathbf n)_k$.
\end{proposition}

\noindent The exact certificate and its verification are given in Appendix~\ref{app:prop:terminal}.

The remaining selected-axis shape has two later CNOTs.  Let wire $j$
meet its second CNOT in chronological position four, on $(j,k)$, and
let the fifth CNOT be $(k,\ell)$ with $j\ne\ell$.  Select an axis on $j$
that commutes through its first incident CNOT.  A collective rotation
can put the fixed final Pauli axis on $\ell$ along $Z_\ell$.  Every such
image lies in the following \emph{overapproximate} space:
\begin{equation}\label{eq:chain-space}
\mathcal L_{jk\ell}=\Span_{\mathbb R}
 \left\{\sigma_\mu^{(j)},\
 \sigma_\mu^{(j)}\sigma_\nu^{(k)},\
 \sigma_\mu^{(j)}\sigma_\nu^{(k)} Z_\ell:
 \mu,\nu\in\{X,Y,Z\}\right\}.
\end{equation}
This space has dimension $3+9+9=21$ and includes either CNOT
orientation and arbitrary intervening local gates.  It is larger than
the actual two-CNOT image family; excluding it is therefore sufficient.

\begin{proposition}[Three-wire certificates]\label{prop:chain}
For the three representatives
\begin{equation}\label{eq:representatives}
(j,k,\ell)=(0,1,2),\quad(0,1,3),\quad(0,2,1),
\end{equation}
the exact Gram matrices of $F$ on \eqref{eq:chain-space} have respective
ranks $207$, $210$, and $189$ among the $231$ quadratic monomials.
The following conclusions hold.
\begin{enumerate}
\item For $(0,1,2)$ and $(0,1,3)$, no nonzero real image in
      $\mathcal L_{jk\ell}$ obeys the spectral-path identity.
\item For $(0,2,1)$, every real image in
      $\mathcal L_{021}$ obeying the identity is supported only on the
      opposite pair $(0,2)$, and its $3\times3$ matrix of pair-Pauli
      coefficients is symmetric.
\end{enumerate}
\end{proposition}

\noindent The exact certificate and its verification are given in Appendix~\ref{app:prop:chain}.

\begin{lemma}[Opposite-first chain excluded]\label{lem:opposite-chain}
The $(0,2,1)$ selected-axis chain of Proposition~\ref{prop:chain}
cannot occur in an exact diagonalizer with five CNOTs.
\end{lemma}

\begin{proof}
The selected image $A$ is a traceless Hermitian involution.  By
Proposition~\ref{prop:chain}, it has a real symmetric pair-Pauli
coefficient matrix $H$.  A collective one-qubit rotation, which
commutes with $V$, diagonalizes $H$ by a proper three-dimensional
rotation.  The transformed image is
\begin{equation}\label{eq:diagonal-pair}
A'=\lambda_X X_0X_2+\lambda_Y Y_0Y_2+\lambda_Z Z_0Z_2.
\end{equation}
The products of distinct terms in \eqref{eq:diagonal-pair} are
$-Z_0Z_2$, $-Y_0Y_2$, or $-X_0X_2$, respectively.  Thus $(A')^2=I$
forces $\lambda_\mu\lambda_\nu=0$ for $\mu\ne\nu$ and
$\sum_\mu\lambda_\mu^2=1$.  Exactly one coefficient is $\pm1$.
Undoing the collective rotation makes
$A=\pm\sigma(\mathbf n)_0\sigma(\mathbf n)_2$.
It commutes with $V^2$, has proper support, and contradicts
Lemma~\ref{lem:trace}.  Dihedral symmetry covers the other
opposite-first representatives.
\end{proof}

\section{From operator constraints to the CNOT bound}

What remains is to show that every short interaction schedule exposes
one of the forbidden observable shapes. At this stage the local gates
and the directions of the CNOTs have already been accounted for. We
only need to classify ordered lists of physical pairs.

There are six unoriented physical pairs.  We enumerate all $6^5=7776$
ordered pair schedules of length five.  This enumeration does not fix
CNOT direction; every test below is necessary for both directions.
For a pair schedule $e_1,\ldots,e_5$, let $d_j$ be the number of
incidences of wire $j$.  A wire's second incident edge is
\emph{terminal} if no later edge meets either endpoint of that edge.

\begin{lemma}[Schedule classification]\label{lem:enumeration}
Every length-five schedule belongs to one of the seven exclusive rows of
Table~\ref{tab:schedules}.  The last three rows each admit a degree-two
wire $j$ whose second incidence is $e_4=(j,k)$ and for which
$e_5=(k,\ell)$ with $j\ne\ell$.  The three rows distinguish the
physical relations among $j,k,\ell$.
\end{lemma}

\begin{table}[ht]
\centering
\caption{Exhaustive ordered, unoriented five-CNOT pair schedules.  Each
row is counted only after the preceding tests fail.}
\label{tab:schedules}
\begin{tabular}{p{0.68\linewidth}r}
\toprule
First necessary obstruction or remaining shape & Count\\
\midrule
Some $d_j<2$ & 5196\\
An input support cone does not reach all four wires & 1572\\
Terminal degree-two edge joins adjacent wires & 576\\
Terminal degree-two edge joins opposite wires & 288\\
$j,k$ adjacent; $j,\ell$ opposite & 48\\
$j,k$ and $j,\ell$ both adjacent & 48\\
$j,k$ opposite & 48\\
\midrule
Total & 7776\\
\bottomrule
\end{tabular}
\end{table}

\begin{proof}
The deterministic classifier in
\texttt{global\_lower\_bound6\_complete\_filter.py} iterates over
the Cartesian product of the six pairs five times.  It first tests
$\min_j d_j<2$, then updates each input support set by adjoining the
endpoints of every edge it touches, then searches for terminal
degree-two edges.  For a schedule passing those tests, it loops over
degree-two wires $j$ in $e_4\setminus e_5$, checks that the other
endpoint $k$ of $e_4$ belongs to $e_5$, and records the third wire
$\ell\in e_5\setminus\{k\}$.  Explicit assertions check that every
schedule reaches exactly one row and that the counts are those shown.
This is a finite exhaustive lemma, independently checkable with integer
and set operations; it imposes no parametrization of local gates.
\end{proof}

\begin{proof}[Proof of Theorem~\ref{thm:main}]
Suppose first that a circuit has exactly five CNOTs.  Its unoriented
pair schedule occurs in Table~\ref{tab:schedules}.  The first row
contradicts Lemma~\ref{lem:degree}, and the second contradicts its
support-cone assertion.  For either terminal-edge row, choose an axis on
the degree-two wire that commutes through its first CNOT.  Because no
later CNOT touches either endpoint of its second edge, its physical
image is a one-CNOT image on that edge.  The adjacent case contradicts
Proposition~\ref{prop:terminal}.  The opposite case is forced by that
proposition to a same-axis opposite-pair product. Since $A^2=I$,
its proportionality factor is $\pm1$, contradicting
Lemma~\ref{lem:trace}.

In each of the last three rows, select the analogous axis on the
degree-two wire $j$.  Its image belongs to the overapproximate space
\eqref{eq:chain-space}.  The first two chain shapes contradict
Proposition~\ref{prop:chain}; the opposite-first shape contradicts
Lemma~\ref{lem:opposite-chain}.  Therefore no five-CNOT diagonalizer
exists.

Fewer than five CNOTs are also excluded.  Lemma~\ref{lem:degree}
already requires at least four.  The following argument handles four.

For completeness, if $m=4$, Lemma~\ref{lem:degree} forces degree two
on every wire.  Choose either endpoint $j$ of the chronological last
CNOT, and select an input axis on $j$ that commutes through its first
incident CNOT.  Its only later CNOT is the last one, which is terminal.
If the last pair is adjacent, its one-CNOT image contradicts
Proposition~\ref{prop:terminal}.  If the pair is opposite, that
proposition forces a same-axis opposite-pair product, contradicting
Lemma~\ref{lem:trace}.  Hence every exact diagonalizer uses at least
six CNOTs.
\end{proof}

\section{Scope of the result}

The obstruction concerns all exact diagonalizers in the stated gate
model, including all-to-all architectures. Restricting connectivity
cannot invalidate it. The theorem leaves a gap between six CNOTs and
the verified fourteen-CNOT construction: it neither constructs a
six-CNOT circuit nor proves the fourteen-CNOT circuit optimal.

The argument has two distinct ingredients. The spectral and trace
identities are derived directly in the text. The small operator-space
exclusions and the schedule classification depend on reproducible exact
computations. No claim of literature novelty is made here; establishing
priority would require a separate literature comparison.

\appendix
\section{Exact spectral-path matrices}\label{app:gram}
For a Pauli basis $E_1,\ldots,E_d$, write
$A(x)=\sum_i x_iE_i$ and let $m(x)$ contain the
$d(d+1)/2$ monomials $x_ix_j$ with $i\leq j$.  Define the nonnegative
spectral-path quartic
\begin{equation}\label{eq:quartic}
F(A)=\sum_{a,b,c\ \mathrm{distinct}}
 \bigl\|M_aAM_bAM_c\bigr\|_{\Frob}^{2}.
\end{equation}
Expand each quadratic matrix $M_aA(x)M_bA(x)M_c$ in the coordinates
$m(x)$ and stack its real and imaginary entries over all $24$ ordered
triples.  If $K$ is the resulting integer matrix, the exact real Gram
matrix is $G=K^{\mathsf T}K$ and
\begin{equation}\label{eq:gram}
F(A(x))=m(x)^{\mathsf T}Gm(x),\qquad
F(A(x))=0\ \Longleftrightarrow\ Gm(x)=0.
\end{equation}
The entries are Gaussian integers before realification.  The Pauli
strings and powers of $V$ have entries in $\{0,\pm1,\pm i\}$, whereas
the projector numerators $M_a$ can have larger entries; each row of
$M_a$ has absolute row sum at most four.  Each quadratic spectral image
in the expansion has absolute row sum at most $128$, and the absolute
value of each Gram entry is bounded by
$24\cdot16^2\cdot128^2<2^{27}$.  The generators use
\texttt{complex128} and check integer-valued output.  These bounds keep
all integer-valued sums in the exact-integer range of binary64.
Nullspaces and Gr\"obner bases are subsequently computed over $\mathbb Q$.


\section{Certificate for a terminal pair}\label{app:prop:terminal}

We verify Proposition~\ref{prop:terminal} as follows.
\begin{proof}
Conjugating a local Pauli axis through a CNOT and arbitrary surrounding
local gates yields a combination of a one-body term on the originating
wire and a two-body term with one fixed Pauli axis on the other wire.
The collective rotation fixes that other axis to $Z$, giving the
overapproximation \eqref{eq:pair-space}.  For the adjacent representative
$(0,1)$, the $21\times21$ Gram matrix from Eqs.~\eqref{eq:projector}, \eqref{eq:quartic}, and
\eqref{eq:gram} has exact rank $21$; it is positive definite.  For the
opposite representative $(0,2)$, its rank is $20$, and its kernel is
spanned by the coordinate $x_{Z_0Z_2}^{2}$.  Because $m(x)$ is made of
actual monomials, every solution in the latter case has only its
$Z_0Z_2$ coefficient nonzero.  Cyclic translations and reflections
transport these facts to every physical pair; reflection exchanges
$V$ with $V^{-1}$ and permutes the labels summed in $F$.
The exact matrices and rank checks are in
\texttt{spectral\_path4\_adjacent\_gram.py} and
\texttt{spectral\_path4\_opposite\_gram.py}.
\end{proof}
\section{Certificates for the three-wire spaces}\label{app:prop:chain}

We verify Proposition~\ref{prop:chain} as follows.
\begin{proof}
In each case, form the basis in Eq.~\eqref{eq:chain-space}, construct
$G$ by \eqref{eq:gram}, and impose $Gm(x)=0$.  A positive-semidefinite
Gram matrix permits exact nullspace elimination of coefficient squares.
For $(0,1,2)$, this forces eight of the $21$ real coefficients to vanish;
for $(0,1,3)$, it likewise forces eight.  With the remaining $13$
variables, append $\sum_i x_i^2-1=0$, permissible by homogeneity for
any nonzero real solution.  In both cases the exact Gr\"obner basis over
$\mathbb Q$ is $\{1\}$, proving item~1.

For $(0,2,1)$, the exact nullspace first forces the three one-body
coefficients and five three-body coefficients to zero.  In the remaining
normalized polynomial ideal, the squares of the four other three-body
coefficients reduce to zero.  Real coefficients therefore make all
three-body terms vanish.  The squares of the antisymmetric differences
of the opposite-pair coefficients also reduce to zero, giving
\begin{equation}\label{eq:symmetric}
x_{X_0Y_2}=x_{Y_0X_2},\qquad
x_{X_0Z_2}=x_{Z_0X_2},\qquad
x_{Y_0Z_2}=x_{Z_0Y_2}.
\end{equation}
These are precisely symmetry of the real pair coefficient matrix.
The stored Gram matrices and rational reductions are generated and
checked by three pairs of scripts with the filename patterns
\path{global_lower_bound6_<case>_gram.py} and
\path{global_lower_bound6_<case>_certificate.py}.
The cases are \path{chain}, \path{adjacent_v}, and
\path{opposite_chain}.
\end{proof}
\section{Reproduction and proof status}

The companion directory contains the deterministic Python scripts and
their saved JSON results.  In particular, run
\texttt{global\_lower\_bound6\_complete\_filter.py} to reproduce
the seven schedule counts and the exact $3\times3$ half-shift trace
tensors $4I_3$ for each opposite pair.  The two
\texttt{spectral\_path4\_*\_gram.py} scripts reproduce the terminal
pair ranks; the six three-wire Gram and certificate scripts reproduce
the three ranks, eliminations, and Gr\"obner reductions.  The main
proof uses exact matrix entries and rational polynomial arithmetic,
not numerical optimization.  Compiling this document confirms only
typesetting; it is separate from checking the mathematical claims and
from independently reviewing the certificate-generating code.
The recorded local verification environment is Python~3.13.9,
NumPy~2.3.5, and SymPy~1.14.0.  The fourteen-CNOT upper bound can be
rechecked with \path{python3 exact_check.py topology14_exact_matchgate_rational.json}.

Theorem~\ref{thm:main} is a lower bound for exact diagonalization in
the specified CNOT plus one-qubit gate model.  It does not assert that
six CNOTs suffice, nor does it give a lower bound for circuits using
ancillas, measurements, other entangling gates, or approximation.

\begin{thebibliography}{9}
\bibitem{cycle-draft}
D.~Miller and A.~A.~Mele, \emph{Exact ancilla-free diagonalization of
qubit cycles}, supplied research draft, September 9, 2026,
Theorem~2 and Sec.~7.
\end{thebibliography}

\end{document}
